\chapter{The grid}\label{ch:grid}
\begin{goals}
\textbf{In this chapter} you learn everything the table itself does: scrolling, sorting, selecting and
copying, widths, number formats, colour scales, hiding and freezing columns, the record inspector and
find.
\end{goals}

The grid shows the current view: the whole file, or what the filter keeps, in the order the sort gives.
Row numbers down the left are the row's position in that view. Above each column header is a small
drawing of the column's distribution, which \cref{ch:gutter} covers.

\section{Scrolling and moving}\label{sec:scrolling}

Scroll with the wheel, the trackpad or the scrollbars, at any depth: only the rows on screen are read
(\cref{sec:screenful}). Type a row number in the row box on the toolbar and press \key{Enter} to jump to
it (\menu{View>Go to row…} puts the cursor in that box).

A trackpad or a trackball rarely scrolls straight, so by default each scroll gesture is locked to the
axis it starts on: start down and you scroll down only. \menu{View>Lock scroll to one axis} turns that
off.

\section{Sorting}\label{sec:sort}

Click a column header to sort by it: ascending, then descending, then back to the file's order. The
header shows an arrow, and the status line says \code{sorted by Cu\_pct ↓}.

To sort by more than one column, right-click the second column's header and choose \ui{Then by …↑} or
\ui{Then by …↓}: the rows keep the first order, and the second breaks its ties. Add as many keys as you
like the same way. \ui{Clear sort} on the same menu, or \menu{Data>Clear sort}, returns to the file's
order.

Numbers sort as numbers and text sorts as text. Blank values sort last, whichever direction. Rows that
tie on every key keep their order in the file.

A sort is a scan (\cref{sec:scans}): on a large file the status line counts it, and \key{Esc} cancels
it. With a filter on, the sort orders only the matches.

\section{Selecting and copying}\label{sec:selection}

Click a cell to select it; drag, or \key{Shift}-click, to select a block of rows and columns. The
selection line in the status bar shows how many cells are selected and, over the numbers among them,
their sum, mean, minimum and maximum: \code{sel 200 · Σ 566.502 · x̄ 2.83251 · min 0.267 · max 7.983}. For
a very large selection it is computed over the first 50{,}000 cells.

\key{Ctrl}+\key{C} copies the selection as tab-separated text, which pastes into a spreadsheet as cells.
The cell's right-click menu also copies it \ui{with header}, \ui{with row \#}, or both.

\section{Column widths}

Drag the border at the right of a column header to resize the column; double-click it to fit the column
to its contents. \menu{View>Autofit all columns} fits every column, and \menu{View>Reset column widths}
goes back to the defaults.

\section{Number formats}\label{sec:formats}

Right-click a number column's header and choose \ui{Number format}: \ui{Auto}, \ui{0 decimals},
\ui{2 decimals}, \ui{3 decimals}, \ui{4 decimals}, \ui{Scientific} or \ui{Thousands (1,234)}. The format
changes only what the grid shows. Filters, statistics and exports use the values as they are in the
file.

\section{Text or number}\label{sec:text-number}

lamina decides from a sample of each column whether it holds numbers or text. When it decides wrong, the
header's menu changes it: \ui{Treat as text} or \ui{Treat as number}. The usual case is a column of codes
that look like numbers (\code{007}, \code{012}): read as numbers they lose their leading zeros and sort
by value. As text they keep them. \Cref{ch:dq} finds such columns for you.

\section{Colour scales}\label{sec:colour-scales}

Right-click a number column's header and choose \ui{Color scale}, then \ui{On}: each cell is painted by
its value. Patterns in a million rows read at a glance this way: the high grades, the empty stretches,
the sentinel values.

\begin{description}
\item[\ui{Palette}] Viridis (the default), Magma, Inferno, Plasma, Turbo, Cividis, Grayscale, and
Blue–red, a diverging palette for values either side of a middle.
\item[\ui{Scale}] Linear or Log. A log scale suits grades, which are usually skewed; it needs positive
values, so a range starting at zero is given a small positive floor.
\item[\ui{Clip outliers (p5–p95)}] stretches the palette over the 5th to 95th percentiles instead of the
whole range, so a few extreme values don't flatten the rest into one colour.
\item[\ui{Reverse}] runs the palette the other way.
\end{description}

Choosing any of these turns the scale on with that setting. The range is taken from the sample behind
the column's distribution (\cref{sec:gutter-sample}), spread through the whole file: its minimum and
maximum, or with clipping, its 5th and 95th percentiles.

\section{Hiding columns}

\ui{Hide …} on a header's menu hides the column; the same menu then lists the hidden columns, to show
them again, and \ui{Show all columns}. Hidden columns are still there: filters and calculated columns
can use them, and an export can include them (\cref{ch:export}).

\section{The columns panel}\label{sec:columns-panel}

\menu{View>Columns…} opens a panel on the right listing every column, for a table with more columns than
fit on the screen.

\shot[0.5\linewidth]{columns-panel}{The columns panel: search, show and hide, the type of each column,
drag handles to reorder, and a pin to freeze a column.}

\begin{itemize}
\item The search box narrows the list as you type.
\item The checkbox beside each column shows or hides it; \ui{Show all}, \ui{Hide all} and \ui{Invert}
act on every column.
\item \ui{\#} or \ui{abc} gives the column's type; a percentage gives the share of blanks in the sample.
\item Drag a column by its \ui{⠿} handle to reorder the grid's columns.
\item Click a column's name to scroll the grid to it.
\item The pin, \ui{\faThumbtack}, \term{freezes} the column: it is moved to the left of the grid and stays there as you
scroll sideways. Freeze the identifiers (\col{ID}, \col{X}, \col{Y}, \col{Z}) of a wide model and they stay
in view. \ui{⋯} opens the column's header menu.
\end{itemize}

Drag the panel's left edge to make it wider; the width is remembered.

\section{The record inspector}\label{sec:inspector}

\menu{View>Record inspector…}, or \ui{Inspect this row…} on a cell's menu, opens a panel listing one row
as name and value, one field per line: the way to read a row with forty columns. It follows the selected
row. \ui{↑} and \ui{↓} step through rows; \ui{⧉} copies the record as tab-separated text.

\ui{⊙} pins the current record. Then step to another and the panel shows the two side by side, the pinned
one and the current one, so you can compare two blocks or two samples field by field.

\section{Find}\label{sec:find}

\key{Ctrl}+\key{F} opens the find bar. Type, and \key{Enter} jumps to the next cell that contains the
text, \key{Shift}+\key{Enter} to the previous one.

\begin{description}
\item[\ui{Aa}] match case.
\item[\ui{⊏⊐}] whole cell: the cell must be exactly the text, not just contain it.
\item[\ui{.*}] a regular expression.
\item[\ui{col}] search only the selected column.
\item[\ui{\#}] count every matching row. That is a scan of the whole view.
\end{description}

Find moves the cursor through the view; it doesn't narrow it. To keep only the rows that match, use the
filter (\cref{ch:filter}): \code{LITO contains "SULPH"}.

\begin{tryit}
Read the sample's richest block.
\begin{enumerate}
\item Open the sample; it is sorted by \col{Cu\_pct}, highest first.
\item Right-click any cell of the first row and choose \ui{Inspect this row…}. The record inspector
lists the block: \col{Cu\_pct} 7.983, \col{LITO} SULPHIDE.
\item Click \ui{⊙} to pin it, then \ui{↓}. The inspector shows the second-richest block beside the first.
\item Right-click the \col{Au\_gpt} header, then \ui{Then by Au\_gpt ↓}: ties on copper, if any, now
put the higher gold first.
\end{enumerate}
\end{tryit}
